\chapter*{Заключение}
\addcontentsline{toc}{chapter}{Заключение}
\label{ch:conclusion}

    В настоящей выпускной квалификационной работе проведено экспериментальное исследование RL-дообучения визуально-языковых моделей действий для повышения эффективности робототехнического манипулирования. По итогам работы все поставленные задачи выполнены в полном объёме.

\section*{Выполнение поставленных задач}
\addcontentsline{toc}{section}{Выполнение поставленных задач}

    \textbf{Задача~1. Аналитический обзор VLA-моделей.} Проведён систематический обзор VLA за 2022--2026~гг. Выделены три архитектурных класса: авторегрессионные (RT-2, OpenVLA), диффузионные ($\pi_0$, FLOWER) и дискретно-диффузионные (dVLA). Зафиксирован экспоненциальный рост публикаций: с 1 работы на ICLR~2024 до 164 на ICLR~2026.

    \textbf{Задача~2. Систематизация подходов интеграции RL и VLM/VLA.} Классифицированы три направления: RL-дообучение VLA (PLD, STARE-VLA -- 98--99\% SR), VLM-генерация наград (EUREKA, Text2Reward) и VLM как планировщик (SayCan, Embodied-R1). Сводная таблица по 8 критериям -- таблица~\ref{tab:integration-classification}.

    \textbf{Задача~3. Разработка методики эксперимента.} Обоснован бенчмарк LIBERO (4 набора, 40 задач). Выбраны девять VLA-моделей, три схемы RL-дообучения, единый бюджет 300~тыс.\ шагов/набор, 3 seed, протокол 10 эпизодов/задачу.

    \textbf{Задача~4. Реализация программного прототипа.} Прототип включает модули VLA-инференса, residual RL, стадийного вознаграждения и VLM-кодогенерации награды. Около 80 запусков (LIBERO + MetaWorld).

    \textbf{Задача~5. Экспериментальное сравнение.} Основные серии~1--4 (VLA+RL, LIBERO) и дополнительная серия~5 (VLM-награды, MetaWorld). Итоги VLA+RL -- таблица~\ref{tab:conclusion_results}; VLM-награды -- таблица~\ref{tab:vlm_metaworld}.

\begin{table}[htbp]
\caption{Итоговые результаты: SR (\%) на LIBERO, среднее по 4 наборам.}
\label{tab:conclusion_results}
\centering
\footnotesize
\setlength{\tabcolsep}{4pt}
\renewcommand{\arraystretch}{1.15}
\begin{tabularx}{\textwidth}{@{}lccc>{\raggedright\arraybackslash}X@{}}
\toprule
\textbf{Модель / метод} & \textbf{До RL} & \textbf{После RL} & \textbf{$\Delta$} & \textbf{Примечание} \\
\midrule
OpenVLA-OFT-7B  & 97{,}1 & 98{,}6 & +1{,}5 & Residual SAC + stage \\
UniVLA-7B       & 95{,}0 & 96{,}8 & +1{,}8 & Residual SAC + stage \\
$\pi_0$-3B      & 94{,}0 & 96{,}1 & +2{,}1 & Residual SAC + stage \\
CogACT-8B       & 93{,}0 & 95{,}4 & +2{,}4 & Residual SAC + stage \\
SmolVLA-2.2B    & 92{,}5 & 95{,}3 & +2{,}8 & Residual SAC + stage \\
$\pi_0$-fast-3B & 90{,}6 & 93{,}8 & +3{,}2 & Residual SAC + stage \\
SpatialVLA-4B   & 88{,}0 & 91{,}5 & +3{,}5 & Residual SAC + stage \\
Octo-Small      & 86{,}0 & 89{,}9 & +3{,}9 & Residual SAC + stage \\
RT-1-X          & 82{,}9 & 87{,}0 & +4{,}1 & Residual SAC + stage \\
\midrule
\multicolumn{5}{@{}l@{}}{\textit{Сравнение RL-схем на OpenVLA-OFT (серия~2)}} \\
Residual SAC + stage  & 97{,}1 & 98{,}6 & +1{,}5 & лучший метод \\
STARE-PPO + stage     & 97{,}1 & 98{,}2 & +1{,}1 & частичная разморозка \\
Direct SAC (unfrozen) & 97{,}1 & 96{,}5 & $-$0{,}6 & деградация \\
\bottomrule
\end{tabularx}
\end{table}

    \textbf{Задача~6. Анализ и рекомендации.} Проведён анализ причин прироста и деградации. Сформулированы пять методических рекомендаций.

\section*{Методические рекомендации}
\addcontentsline{toc}{section}{Методические рекомендации}

    По итогам экспериментов сложилось несколько практических рекомендаций. Прежде всего, выбирать VLA имеет смысл с оглядкой на её исходный SR. У сильных моделей вроде OpenVLA-OFT (97,1\%) RL добавляет немного, +1--2~п.п., но уверенно выводит систему за 98\%; у моделей среднего уровня (SmolVLA, $\pi_0$, Octo, 83--93\%) прирост заметнее, +3--4~п.п.; RT-1-X (82,9\%) выигрывает от дообучения больше всех (+4,1~п.п.), хотя по абсолютному SR всё равно остаётся позади OpenVLA.

    Residual-схему с замороженной VLA по нашим данным стоит предпочесть полной разморозке. Полная разморозка (direct SAC) теряет 0,6~п.п.\ из-за катастрофического забывания, тогда как лёгкая надстройка (MLP 256/256 на эмбеддингах VLA и проприоцепции) сохраняет накопленные навыки и лишь подправляет остаточные ошибки. Эта рекомендация получена на ограниченной выборке (3 seed, 10 эпизодов на задачу), поэтому её стоит воспринимать как практическое наблюдение, а не как строго доказанное утверждение. Для длинных задач при этом важна стадийность награды: разбиение на Reach $\to$ Grasp $\to$ Transport $\to$ Place даёт +0,8~п.п.\ над разреженным сигналом на Spatial и до +2,5~п.п.\ на наборе Long.

    В качестве алгоритма для residual RL лучше брать SAC, а не PPO: его буфер опыта сглаживает нестабильность стадийных сигналов и в итоге даёт +1,5 против +1,1~п.п.\ у STARE-PPO (PPO остаётся разумным выбором при частичной разморозке, но требует аккуратной настройки). Наконец, само направление интеграции стоит выбирать по наличию готовой VLA. Если есть модель с SR\,$>$\,90\%, имеет смысл идти через residual SAC со стадийной наградой на LIBERO; если же предобученной VLA нет, остаётся путь через VLM-генерацию награды с последующим RL, где на контактных задачах помогает 3--5 итераций рефлексии (в серии~5 push поднялся с 48 до 85\%).

\section*{Основные результаты и научная значимость}
\addcontentsline{toc}{section}{Основные результаты и научная значимость}

    Главный результат работы в том, что residual RL со стадийной наградой устойчиво поднимает SR всех девяти рассмотренных VLA: прирост на LIBERO составил от +1,5 до +4,1~п.п., а OpenVLA-OFT после дообучения вышла на 98,6\% -- вплотную к 99\% из PLD~\cite{pld2026}, но при меньшем бюджете. Полная же разморозка модели в наших экспериментах оказалась вредной: direct SAC снижает SR с 97,1 до 96,5\%, поэтому заморозку базовой VLA мы рекомендуем как способ стабилизировать обучение (с поправкой на ограниченный объём проведённых экспериментов).

    Любопытна и обнаруженная закономерность: чем слабее исходная модель, тем больше она выигрывает в абсолютных числах (RT-1-X), а сильные модели прибавляют меньше, зато достигают более высокого итогового SR. Формально это выражается сильной обратной корреляцией ($\rho = -0{,}97$) между baseline SR и приростом от RL. Стоит, впрочем, оговориться, что часть этой зависимости -- следствие близости сильных моделей к потолку бенчмарка (запас для роста у них мал по определению, не более 1--3~п.п.), поэтому эффект не специфичен для RL, и мы рассматриваем его как практическое наблюдение, а не как фундаментальное свойство метода. Подтвердилась и роль стадийности награды: без неё прирост падает с +1,5 до +0,7~п.п., поскольку именно стадии позволяют разложить длинные траектории LIBERO-Long. Все девять моделей -- OpenVLA-OFT, UniVLA, $\pi_0$, CogACT, SmolVLA, $\pi_0$-fast, SpatialVLA, Octo и RT-1-X -- сравнивались в едином протоколе, что отличает работу от большинства статей, ограничивающихся одной-двумя моделями.

    Дополнительная серия показала, что продвинуться можно и без предобученной VLA: итеративная VLM-кодогенерация награды поднимает средний SR на MetaWorld с 53,9 до 77,2\% за три итерации, причём корреляция $\rho(R,S) > 0{,}8$ служит надёжным предиктором успеха, а SAC и здесь обходит PPO (reach: 100 против 71,7\%). В сумме результат 98,6\% при 300k шагов на набор оказывается всего на 0,4~п.п.\ ниже PLD (99,0\% при $\geq$1M шагов), то есть остаётся конкурентоспособным при существенно меньших затратах.


\section*{Направления дальнейших исследований}
\addcontentsline{toc}{section}{Направления дальнейших исследований}

    Работу естественно продолжить в нескольких направлениях. Самое прямое -- увеличить бюджет дообучения: при 500k--1M шагов на набор OpenVLA, по аналогии с PLD~\cite{pld2026}, должна вплотную подойти к 99\% SR. Дальше напрашивается дистилляция: собранные RL-траектории можно влить обратно в VLA через supervised fine-tuning и тем самым убрать residual-надстройку на этапе инференса. Отдельная большая задача -- проверка sim-to-real, то есть прогон дообученных моделей на реальном манипуляторе Franka и оценка того, как разрыв между симуляцией и реальностью влияет на полученный прирост. Наконец, оба трека работы можно объединить в единый конвейер: VLM-награда для сбора данных, затем дистилляция в VLA и, наконец, residual RL для финальной коррекции.

\section*{Заключительные положения}
\addcontentsline{toc}{section}{Заключительные положения}

    \textit{Теоретический вклад}: систематическая классификация более 20 методов интеграции RL и VLM/VLA (включая свежие работы 2025--2026~гг.\ по RL-дообучению VLA: VLA-RL, SimpleVLA-RL, RIPT-VLA, ConRFT, VLA-RFT, RLinf-VLA); обоснование выбора residual RL со стадийным вознаграждением и позиционирование работы относительно полнополитических токен-уровневых подходов (PPO/GRPO/LOOP).

    \textit{Экспериментальный вклад}: основной блок -- 9 VLA и 3 RL-схемы на LIBERO; дополнительный -- VLM-награды на MetaWorld (около 80 запусков суммарно).

    \textit{Практический вклад}: пять рекомендаций по выбору VLA, схемы RL-дообучения и протокола оценки; воспроизводимый прототип на Python.

    Результаты подтверждают: RL-дообучение VLA -- практичный путь к SR\,$>$\,98\% при наличии предобученной модели с SR\,$>$\,90\%.

\endinput
